% problem_id: S051-drops-D4
% source_id: S051 (DROPS/LIPIcs)
% source_file: t30.pdf
% statement_locator: printed p. 30:10 (PDF page 10); the numbered items 1-3 of the statement continue onto 30:11 (PDF page 11). Prerequisite Lemma 11 is stated on 30:10, and Definition 9 (active/passive points) on 30:9.
% proof_locator: printed p. 30:11 (PDF page 11) for Lemma 12; 30:10 (PDF page 10) for the proof of the prerequisite Lemma 11 (and for Lemma 10, the proof that ray shootings only start at p_a or at active points). Context pages read verbatim for the chain: 30:9 (PDF 9, module 3.2 opening: naive ray shooting + Definition 9), 30:4 (PDF 4, Theorem 4), 30:6 (PDF 6, Theorem 7), 30:14 (PDF 14, Lemma 14, downstream).
% representation: PDF; transcribed from rendered page images (never from the text layer)
% collected_by: carrier rule (the headline theorem was an assembly)
% review_status: reviewers split 1-1 => adjudicated H2 (burden of proof on H3)
% status: complete (statement + author proof verbatim + reason + locators)
%
% =====================================================================
% Candidate: D4
% Paper: Itai Boneh, Shay Golan, Shay Mozes, Oren Weimann,
%        "Optimal Dynamic Time Warping on Run-Length Encoded Strings",
%        LIPIcs Vol. 297, ICALP 2024, Article No. 30, pp. 30:1-30:17.
%        DOI 10.4230/LIPIcs.ICALP.2024.30   (full version: arXiv:2302.06252, ref [9])
% Source file: /disks/sata1/yupeng/human-proof-corpus/source-cache/registry-sources/drops/t30.pdf
% Rendered with: pdftoppm -f <p> -l <p> -r 150 -png   (pdftotext used ONLY for locating)
% Page mapping: PDF page k == printed page 30:k.
%
% COLLECTED UNIT (see D4.json):
%   Lemma 12 (long/short ray dichotomy), printed pp. 30:10-30:11,
%   with its immediate prerequisite Lemma 11 (printed p. 30:10),
%   within the module 3.2 "Active and Passive Points, Long and Short Rays"
%   (Definition 9, Lemma 10, mega-segments, Lemma 11, Lemma 12).
%   Transcribed here: full printed pages 30:9, 30:10, 30:11.
%
% Verbatim transcription. No rewriting, no completion. Original typos are
% kept and flagged. Illegible spots are flagged as %% [ILLEGIBLE: ...].
% =====================================================================

\documentclass[11pt]{article}
\usepackage{amsmath}
\usepackage{amssymb}
\usepackage[margin=1in]{geometry}
\usepackage{enumitem}
\newcommand{\Pset}{\mathcal{P}}
\newcommand{\Pact}{\mathcal{P}_{\mathrm{active}}}
\newcommand{\Lwave}{\ensuremath{\mathsf{LeftLinearWave}}}
\newcommand{\Rwave}{\ensuremath{\mathsf{RightLinearWave}}}
\newcommand{\AddConst}{\ensuremath{\mathsf{AddConst}}}
\newcommand{\AddGrad}{\ensuremath{\mathsf{AddGradient}}}
\newcommand{\nextGT}{\ensuremath{\mathsf{nextGT}}}
\newcommand{\prevLT}{\ensuremath{\mathsf{prevLT}}}
\newcommand{\poly}{\ensuremath{\mathrm{polylog}}}
\newcommand{\Fset}{\ensuremath{\mathsf{F}}}
\begin{document}

\section*{D4 --- Transcription (verbatim)}

% ---------------------------------------------------------------------
% CONTEXT A (verbatim, printed p. 30:4): the result the paper advertises
% as its "main technical contribution", quoted so that the reduction
% chain can be checked. Its complete proof is NOT in the conference paper.
% ---------------------------------------------------------------------
\subsection*{Context A --- Theorem 4 (printed p.~30:4, PDF p.~4)}

In Section 3, we show our main technical contribution:

\smallskip
$\blacktriangleright$ \textbf{Theorem 4.} \emph{Performing $s$ range operations of Definition 3 can be done in amortized $O(\poly(s))$ time per operation.}

\smallskip
%% [Editorial note, not transcription: Section 3 (printed pp. 30:6-30:14) opens with
%%  "In this section, we prove Theorem 4." but section 3.3 (printed p. 30:12) states that
%%  what is given is a "technical overview" and that "The complete implementation details
%%  and proofs appear in the full version of this paper ([9])." Hence Theorem 4 has no
%%  complete proof inside this conference paper.]

% ---------------------------------------------------------------------
% CONTEXT B (verbatim, printed p. 30:6): the flagship theorem. Its proof is a
% two-paragraph reduction (printed p. 30:6), quoted in D4.json only in summary.
% ---------------------------------------------------------------------
\subsection*{Context B --- Theorem 7 (printed p.~30:6, PDF p.~6), the flagship}

$\blacktriangleright$ \textbf{Theorem 7.} \emph{The Dynamic Time Warping distance of two run-length encoded strings $S$ and $T$ with $n$ and $m$ runs respectively can be computed in $\tilde{O}(nm)$ time.}

\medskip
%% [Editorial note, not transcription: the proof of Theorem 7 (printed p. 30:6) is a
%%  reduction only: it initialises the data structure of Theorem 4, runs $nm$ iterations,
%%  and applies "$O(1)$ range operations (due to Lemma 6)" per iteration. Lemma 6
%%  (printed p. 30:5: "$\Fset_{t+1}$ can be obtained by using $O(1)$ range operations
%%  (Definition 3) on $\Fset_t$.") is proved, per the authors, only in the full version
%%  ("In the full version of this paper ([9]), we prove the following lemma.").]

\bigskip
% =====================================================================
% THE COLLECTED UNIT: printed pages 30:9, 30:10, 30:11 (PDF pages 9, 10, 11)
% =====================================================================

%% --- page 9 (printed 30:9; running head: I. Boneh, S. Golan, S. Mozes, and O. Weimann) ---

$\alpha$ is shot from the next $p_w$ with $\alpha_w > \alpha$, and so on. In the the full version of this paper
([9]) we formally prove that the above characterization indeed represents the new values of
$A[i \ldots j]$.
%% [ORIGINAL TYPO: "the the" (doubled article) in the original; verified at 400 dpi.]

We now describe a naive, non-efficient implementation of \Lwave$(i,j,\alpha)$ according
to the description above. Recall that $i$ and $j$ are assumed to be endpoints $p_a$ and $p_{b+1}$
of segments. We begin by finding the first $p_z$ with $x_z \in [i \ldots j]$ and $\alpha_z > \alpha$ by querying
$D_\alpha.\nextGT(i,\alpha)$. A ray shooting process is then performed from $p_z$ (if $p_z$ exists) as follows:
Recall that $r_z(x)$ denotes the positive ray with slope $\alpha$ shot from $p_z$. We scan the successor
points of $p_z$ one by one in order, and for every point $p_w$ we check whether the ray $r_z(x_w) \le y_w$.
If so, $p_w$ is removed by removing $x_w$ from $\Pset$, $D_\beta$ and $D_\alpha$. Otherwise, we compute $p^* = (x^*, y^*)$,
the intersection point of $r_z$ and $\ell_{w-1}$, and from it the points $p' = (x', y') = (\lfloor x^* \rfloor, r_z(\lfloor x^* \rfloor))$
and, if $x^*$ is not an integer, also $p'' = (\lceil x^* \rceil, \ell_{w-1}(\lceil x^* \rceil))$. Then, we insert the new points $p'$
and $p''$ just before $p_w$, as discussed above for \AddConst. The scanning then continues with
another \nextGT{} query from $p_w$, and so on. If, at the end of the process, the last point $p_b$ is
removed since it is above some $r_z$, we insert a new point $(x_b, r_z(x_b))$.

\smallskip
\noindent\textbf{Time Complexity.}\quad We now analyze the time complexity of this naive implementation.
Each \AddConst{} and \AddGrad{} operation requires $O(1)$ operations on the Interval-add
data structures, and therefore takes $O(\poly|\Pset|)$ time per operation, with $|\Pset|$ being the
cardinality of $\Pset$ when the operation is applied.

Regarding \Lwave{} operations, one might hope that the cost of each ray shooting
process can be charged to the removal of points from $\Pset$ during the process. However, each
ray shooting process might also add up to two new points, which might result in the size of
$\Pset$ increasing. Indeed, a \Lwave{} operation may give rise to many such ray shooting
processes, and hence may significantly increase the size of $\Pset$ and take too much time. This
is the main technical challenge we need to address.

The idea is to distinguish between \emph{long} ray shootings for which we can globally charge
the new insertions, and \emph{short} ray shootings for which we cannot. We handle the long rays as
in the naive solution and devise a separate lazy mechanism that delays the application of all
the short rays stemming from a single \Lwave{} operation using a constant number of
updates to a separate data structure that keeps track of the delayed rays.

\smallskip
\noindent\textbf{Symmetry of \Rwave.}\quad The discussion so far was focused on the \Lwave{}
operation. We note that the analysis of \Rwave{} is symmetric. In particular, the
execution of \Rwave$(i,j,\alpha)$ can be described as a sequence of ray shootings with
\emph{negative} rays. The first point from which a ray is shot is $p_z$ with largest $z \in [a \ldots b]$ such that
$\alpha_{z-1} < -\alpha$ ($p_z$ is found using $D_\alpha.\prevLT$). Note that the condition for starting a ray shooting
process for \Rwave{} is on $\alpha_{z-1}$ rather than $\alpha_z$ since the slope of the segment to the
left of $p_z$ is $\alpha_{z-1}$. To simplify the presentation, we will keep describing only \Lwave,
and will comment at the very end about the minor adjustments required to also handle the
symmetric \Rwave.

\subsection*{3.2\quad Active and Passive Points, Long and Short Rays}

On our way to formally define long rays and short rays we first observe that ray shootings
only occur at points where slopes increase. We call such points \emph{active} points.

\smallskip
$\blacktriangleright$ \textbf{Definition 9 (Active and Passive points).} \emph{A point $p_z$ in $\Pset$ is called active if $z \in \{1, |\Pset|\}$ or
$\alpha_z > \alpha_{z-1}$. A point that is not active, is called passive. We denote the sets of active points
by $\Pact$.}

%% [footer of printed p. 30:9: "ICALP 2024"]

%% --- page 10 (printed 30:10) ---

$\blacktriangleright$ \textbf{Lemma 10.} \emph{Ray shootings stemming from \Lwave$(i,j,\alpha)$ occur either at point
$p_a = (i, A[i])$ or at active points. Ray shootings stemming from \Rwave$(i,j,\alpha)$ occur
either at point $p_b = (j, A[j])$ or at active points.}

\smallskip
\noindent\textbf{Proof.}\quad We focus on \Lwave. The proof for \Rwave{} is symmetric. Assume
to the contrary that a ray shooting process starts at a passive point $p_z \neq p_a$. If $p_z$ is the
first point where a ray shooting starts, then $z$ is the minimal index in $[a \ldots b]$ with $\alpha_z > \alpha$.
But since $p_z$ is passive, we have $\alpha < \alpha_z \le \alpha_{z-1}$, contradicting the minimality of $z$ (note that
$p_z \neq p_a$ so $z - 1 \in [a \ldots b]$).

Otherwise, let $p_q$ be the last point before $p_z$ from which a ray shooting process occurred.
Let $p_{q'}$ be the first point below the ray shot from $p_q$. Since $p_z$ is the next point from which
a ray is shot, $z$ is the first point in $[q' \ldots b]$ with $\alpha_z \ge \alpha$. Since $p_z$ is passive, we have
$\alpha < \alpha_z \le \alpha_{z-1}$. If $z \neq q'$, we have $z - 1 \in [q' \ldots b]$, a contradiction to the minimality of $z$.
Otherwise, $p_z = p_{q'}$ is the first point below the ray with slope $\alpha$ shot from $p_q$. It follows
that $p_{z-1}$ is above the ray, and $\alpha_{z-1} > \alpha$. It must be the case that $p_{q'}$ is above the ray, a
contradiction. \hfill$\blacktriangleleft$

\smallskip
Let $\Pact = (q_1, q_2, \ldots)$ be the restriction of the sequence $\Pset$ to the active points. We
can think of the active points as defining a piecewise linear function whose segments are
a coarsening of the segments of $A$. We refer to these segments as \emph{mega-segments}. Let $\gamma_z$
denote the slope of the mega-segment whose endpoints are $q_z$ and $q_{z+1}$. The following lemma
asserts that the segments of $A$ are never below their corresponding mega-segments, and that
the slope of a segment starting at an active point is never smaller than the slope of the
mega-segment starting at the same point.

\smallskip
$\blacktriangleright$ \textbf{Lemma 11.} \emph{Let $q_z = p_w$ and $q_{z+1} = p_{w'}$ be two consecutive active points. For every
$k \in [w \ldots w']$, the passive point $p_k$ is not below the mega-segment connecting $q_z$ and $q_{z+1}$.
Furthermore, $\alpha_w \ge \gamma_z$.}

\smallskip
\noindent\textbf{Proof.}\quad (See Figure 4) Clearly, $p_w$ and $p_{w'}$ are on the mega-segment, and in particular not
below it. Assume by contradiction that there is a point below the mega-segment, and let
$k' \in (w \ldots w')$ be the smallest index of such a point. Since $p_{k'-1}$ is not below the mega-
segment and $p_{k'}$ is below the mega-segment, we must have $\alpha_{k'-1} < \gamma_z$. Moreover, since the
points $p_k$ with $k \in [k' \ldots w')$ are passive, the slopes are non-increasing and therefore every
$\alpha_k \le \gamma_z$. This means that all these points and in particular $p_{w'}$ are below the mega-segment.
In contradiction to $p_{w'}$ lying on the mega-segment. Furthermore, since $p_{w+1}$ is not below
the mega-segment, we have $\alpha_w \ge \gamma_z$. \hfill$\blacktriangleleft$

\smallskip
We next show that if a ray shooting process starts at an active point $q_z$ with $\gamma_z < \alpha$ then
the process ends before $q_{z+1}$, and the only affected points are the passive points between $q_z$
and $q_{z+1}$. On the other hand, if $\gamma_z \ge \alpha$ then as a result of the process $q_{z+1}$ ceases to be an
active point, so $|\Pact|$ decreases.

\smallskip
$\blacktriangleright$ \textbf{Lemma 12.} \emph{Consider a ray shooting process starting from point $p_w = q_z \in \Pact$ during
the application of \Lwave$(i,j,\alpha)$. Let $q_{z+1} = p_{w'}$.}
\begin{enumerate}[label=\arabic*., leftmargin=*, itemsep=0pt, topsep=2pt]
\item \emph{No new active points $p = (x, y)$ with $x \neq j$ are created in this process.}
\item \emph{If $\gamma_z < \alpha$ then the points that are deleted by this process are the (passive) points $p_k$ with
$k \in [w + 1 \ldots r]$ for some $w < r < w'$. No other points between $p_w$ and $p_{w'}$ are deleted by
\Lwave$(i,j,\alpha)$.}
\item \emph{If $\gamma_z \ge \alpha$ then $q_{z+1}$ is either deleted or becomes passive.}
\end{enumerate}

%% --- page 11 (printed 30:11; running head: I. Boneh, S. Golan, S. Mozes, and O. Weimann) ---

%% [Figure 4 occupies the top of printed p. 30:11. Caption, verbatim:]
\begin{quote}
\small
Figure 4\quad The impossible configuration in Lemma 11. The points $q_z$ and $q_{z+1}$ are represented by
the purple points and the mega-segment connecting them is represented by a thick purple line. The
first point $p_{k'}$ below the segment is marked with red stroke. Since the points strictly within the
mega-segment are passive, the points following $p_{k'}$ within the mega-segment (and in particular $q_{z+1}$)
must remain below the mega-segment.
\end{quote}

\noindent\textbf{Proof.}\quad Let $\ell$ be the ray starting from $p_w = q_z$. Assume the process terminates by finding the
first point $p^* = (x^*, y^*)$ below $\ell$ (the only process that does not end this way is the one that
ends by reaching $(j, A[j])$). The ray shooting process adds at most two new points $p'$ and $p''$
with decreasing slopes, so no new active points are created by the process. The slope of $p'$ is
decreasing because the segment entering $p'$ is (a sub-segment of) $\ell$ and the segment leaving
$p'$ is to a point below $\ell$. The slope of $p''$ is decreasing because the line segment entering $p''$ is
a line from $p'$ (a point on $\ell$) and the line segment leaving $p''$ is to the suffix of a line segment
below $\ell$.

Consider the case $\gamma_z < \alpha$. Then $q_{z+1}$ is below the ray with slope $\alpha$ starting at $q_z$. Hence
the ray shooting process terminates at a point after $p_{w+1}$ and before $q_{z+1}$. Since no active
points are created, the next ray will be shot from $q_{z+1}$ or later, so no other points between
$q_z$ and $q_{z+1}$ are deleted by \Lwave$(i,j,\alpha)$.

Now consider the case $\gamma_z > \alpha$. Then the mega-segment between $q_z$ and $q_{z+1}$ is above the
ray with slope $\alpha$ shot from $q_z$. By Lemma 11, all the (passive) points between $q_z$ and $q_{z+1}$
are also above this ray. Hence $q_{z+1}$ is deleted by the ray shooting process.

Finally, consider the case $\gamma_z = \alpha$. Then the mega-segment between $q_z$ and $q_{z+1}$ coincides
with the ray with slope $\alpha$ shot from $q_z$. By Lemma 11, all the (passive) points between $q_z$
and $q_{z+1}$ will be deleted by the ray shooting process. Let $w'$ be such that $q_{z+1} = p_{w'}$. If
$\alpha_{w'} \ge \alpha$ then $q_{z+1}$ will be deleted by the process. Otherwise, $\alpha_{w'} < \alpha$, so the ray shooting
process terminates at $q_{z+1}$. Since all the passive points between $q_z$ and $q_{z+1}$ were deleted, $q_z$
and $q_{z+1}$ become consecutive in $\Pset$, and the slope of the corresponding segment is $\gamma_z = \alpha$.
But the slope of the segment starting at $q_{z+1}$ is $\alpha_{w'} < \alpha$, so $q_{z+1}$ becomes passive. \hfill$\blacktriangleleft$

\smallskip
We call rays with $\gamma_z > \alpha$ \emph{long} rays, and those with $\gamma_z \le \alpha$ \emph{short} rays. Since long rays
decrease the size of $\Pact$ we can handle them explicitly as in the warmup, charging the
deletion of passive points during the process to their creation, and charging the insertion
of the at most two passive points at the end of the process to the decrease in $|\Pact|$. The
short rays, which do not decrease $|\Pact|$, will be handled lazily. Namely, instead of explicitly
shooting a short ray in the mega-segment starting at an active point $q_z$, we only store the
slope of the ray and postpone its execution until it is required (e.g., by a \ensuremath{\mathsf{Lookup}} operation).
Note that subsequent short rays shot in this mega-segment may further change the stored
slope, and subsequent long rays may also affect it. We explain this in detail next.

%% [footer of printed p. 30:11: "ICALP 2024"]

\bigskip
%% ---------------------------------------------------------------------
%% DOWNSTREAM (not transcribed here, listed for the judgment in D4.json):
%%  * printed p. 30:12: "3.3 The Data Structure" -- Invariants 1-3
%%    (Ptilde_active = P_active; rho_z > gamma_z; A[x] = min(ell_w(x), r_z(x))),
%%    and the explicit statement that Section 3 is a "technical overview" whose
%%    "complete implementation details and proofs appear in the full version".
%%  * printed p. 30:12: Lemma 13 (Add-min Data Structure) -- proof in the full version.
%%  * printed p. 30:13-30:14: the flush operation.
%%  * printed p. 30:14: Lemma 14 ("Applying flush to an active point q_z in P_active
%%    preserves Invariants 1-3. ...") -- complete proof in the conference paper;
%%    its first step uses Lemma 12 ("it is guaranteed by Lemma 12 that the scan of
%%    flush ends at q_{z+1} or before q_{z+1}").
%%  * printed p. 30:14-30:15: Claim 16 + Theorem 15 (bounded DTW) -- complete proofs,
%%    routine k-band counting argument.
%% ---------------------------------------------------------------------
\end{document}
